\subsection{DUAL OF A VECTOR SPACE}

THEOREM 4. The dimension of the dual \(\mathbf{E}^*\) of a vector space \(\mathbf{E}\) is at least equal to the dimension of \(\mathbf{E}\) . For \(\mathbf{E}^*\) to be finite-dimensional, it is necessary and sufficient that \(\mathbf{E}\) be so, and then \(\dim \mathbf{E}^* = \dim \mathbf{E}\) .

If K is the field of scalars of E, E is isomorphic to a space \(\mathbf{K}_{s}^{\mathrm{(I)}}\) and therefore \(E^{*}\) is isomorphic to \(K_{d}^{I}\) (§2, no. 6, Proposition 10). As \(\mathbf{K}_{s}^{\mathrm{(I)}}\) is a subspace of \(K_{d}^{I}\) , \(\dim E = \text{Card}(I) \leqslant \dim E^{*}\) (no. 3, Corollary 4 to Proposition 4); further, if I is finite, \(\mathbf{K}_{d}^{I} = \mathbf{K}_{d}^{\mathrm{(I)}}\) (cf.~Exercise 3(d)).

COROLLARY. For a vector space \(\mathbf{E}\) , the relations \(\mathbf{E} = \{0\}\) and \(\mathbf{E}^{*} = \{0\}\) are equivalent.

THEOREM 5. Given two exact sequences of vector spaces (over the same field K) and linear mappings

\[
0 \rightarrow \mathrm{E} ^ {\prime} \rightarrow \mathrm{E} \rightarrow \mathrm{E} ^ {\prime \prime} \rightarrow 0
\]

\[
0 \to \mathbf {F} ^ {\prime} \to \mathbf {F} \to \mathbf {F} ^ {\prime \prime} \to 0
\]

and two vector spaces G, H over K, the corresponding sequences

\[
0 \rightarrow \operatorname{Hom} (\mathrm{E} ^ {\prime \prime}, \mathrm{G}) \rightarrow \operatorname{Hom} (\mathrm{E}, \mathrm{G}) \rightarrow \operatorname{Hom} (\mathrm{E} ^ {\prime}, \mathrm{G}) \rightarrow 0
\]

\[
0 \rightarrow \operatorname{Hom} (\mathrm{H}, \mathrm{F} ^ {\prime}) \rightarrow \operatorname{Hom} (\mathrm{H}, \mathrm{F}) \rightarrow \operatorname{Hom} (\mathrm{H}, \mathrm{F} ^ {\prime \prime}) \rightarrow 0
\]

are exact and split.

This follows from the fact that every vector subspace is a direct factor (no. 3, Proposition 4) and from § 2, no. 1, Propositions 1 and 2.

COROLLARY. For every exact sequence

\[
0 \longrightarrow \mathrm{E} ^ {\prime} \stackrel {u} {\longrightarrow} \mathrm{E} \stackrel {v} {\longrightarrow} \mathrm{E} ^ {\prime \prime} \longrightarrow 0
\]

of vector spaces over the same field K and linear mappings the sequence

\[
0 \longrightarrow \mathrm{E} ^ {\prime \prime *} \stackrel {t _ {v}} {\longrightarrow} \mathrm{E} ^ {*} \stackrel {t _ {u}} {\longrightarrow} \mathrm{E} ^ {\prime *} \longrightarrow 0
\]

is exact and splits.

It follows in particular that for every vector subspace M of E the canonical homomorphism \(E^{*}/M^{\prime} \rightarrow M^{*}\) , where \(M^{\prime}\) is the subspace of \(E^{*}\) orthogonal to M ( \(\S 2\) , no. 4), is bijective.

THEOREM 6. For every vector space E over a field K, the canonical mapping \(c_{\mathbf{E}}: \mathbf{E} \to \mathbf{E}^{**}\) (§2, no. 7) is injective; for it to be bijective, it is necessary and sufficient that E be finite-dimensional.

The first assertion and the fact that if E is finite-dimensional \(c_{E}\) is bijective are special cases of §2, no. 7, Proposition 14. Suppose that E is infinite-dimensional, so that we may assume that \(\mathbf{E} = \mathbf{K}_{s}^{(L)}\) , where L is an infinite set and therefore \(E^{*} = K_{d}^{L}\) . Let \((e_{\lambda})_{\lambda \in L}\) be the canonical basis of E and let \((e_{\lambda}^{*})_{\lambda \in L}\) be the corresponding family of coordinate forms in \(E^{*}\) (§2, no. 6); the vector subspace of \(E^{*}\) generated by the \(e_{\lambda}^{*}\) is just the direct sum \(F' = K_{d}^{(L)}\) and the hypothesis that L is infinite implies that \(F' \neq E^{*}\) . Then there exists a hyperplane \(H'\) of \(E^{*}\) containing \(F'\) (no. 3, Proposition 8) and, as \(E^{*}/H'\) is non-zero, its dual is also non-zero (Corollary to Theorem 4), which is identified with the orthogonal \(H''\) of \(H'\) in \(E^{**}\) (§2, no. 6, Corollary to Proposition 9). But \(H'' \cap c_{E}(E)\) is contained in the image under \(c_{E}\) of the orthogonal of \(F'\) in E, which is by definition 0; hence \(c_{E}(E) = E^{**}\) is impossible.

E will usually be identified with the subspace of \(E^{**}\) the image of \(c_{E}\) .

Let \(\mathbf{E}, \mathbf{F}\) be two vector spaces over a field \(\mathbf{K}\) and \(u: \mathbf{E} \to \mathbf{F}\) a linear mapping. We shall define canonical isomorphisms:

\begin{enumerate}
\def\labelenumi{(\arabic{enumi})}
\item
  Of the dual of \(\operatorname{Im}(u) = u(\mathbf{E})\) onto \(\operatorname{Im}(^t u) = {}^t u(\mathbf{F}^*)\) .
\item
  Of the dual of \(\operatorname{Ker}(u) = \overline{u}^1(0)\) onto \(\operatorname{Coker}(^t u) = \mathrm{E}^* / {}^t u(\mathrm{F}^*)\) .
\item
  Of the dual of \(\operatorname{Coker}(u) = \mathbf{F} / u(\mathbf{E})\) onto \(\operatorname{Ker}(^t u) = {}^t u^{-1}(0)\) .
\end{enumerate}

We write \(\mathbf{I} = \operatorname{Im}(u),\mathbf{N} = \operatorname{Ker}(u),\mathbf{C} = \operatorname{Coker}(u)\) ; from the exact sequences

\[
0 \longrightarrow \mathrm{N} \longrightarrow \mathrm{E} \stackrel {p} {\longrightarrow} \mathrm{I} \longrightarrow 0, \quad 0 \longrightarrow \mathrm{I} \stackrel {j} {\longrightarrow} \mathrm{F} \longrightarrow \mathrm{C} \longrightarrow 0 \tag {14}
\]

we derive, by transposition (Corollary to Theorem 5), the exact sequences

\[
\begin{array}{l} 0 \longrightarrow \mathrm{I} ^ {*} \stackrel {{t _ {p}}} {{\longrightarrow}} \mathrm{E} ^ {*} \longrightarrow \mathrm{N} ^ {*} \longrightarrow 0, \\ 0 \longrightarrow \mathrm{C} ^ {*} \longrightarrow \mathrm{F} ^ {*} \stackrel {{t _ {j}}} {{\longrightarrow}} \mathrm{I} ^ {*} \longrightarrow 0. \end{array}\tag{15}
\]

Moreover, as \(u = j \circ p\) , \(^{t}u = ^{t}p \circ ^{t}j\) ; the exact sequences (15) thus define canonical isomorphisms of \(C^*\) onto \(\operatorname{Ker}(^{t}u)\) , of \(I^*\) onto \(\operatorname{Im}(^{t}u)\) and of \(N^*\) onto \(\operatorname{Coker}(^{t}u)\) , since \(^{t}p\) is injective and \(^{t}j\) surjective. To be precise, let \(y \in \operatorname{Im}(u)\) , \(z \in \operatorname{Ker}(u)\) , \(t \in \operatorname{Coker}(u)\) , \(y' \in \operatorname{Im}(^{t}u)\) , \(z' \in \operatorname{Coker}(^{t}u)\) , \(t' \in \operatorname{Ker}(^{t}u)\) ; when \(y'\) , \(z'\) , \(t'\) are canonically identified with linear forms on \(\operatorname{Im}(u)\) , \(\operatorname{Ker}(u)\) and \(\operatorname{Coker}(u)\) respectively, then

\[
\langle y, y ^ {\prime} \rangle = \langle x, y ^ {\prime} \rangle \quad \text { for   all } x \in \mathrm{E} \text { such   that } u (x) = y;\tag{16}
\]

\begin{enumerate}
\def\labelenumi{(\arabic{enumi})}
\setcounter{enumi}{16}
\item
  \(\langle z, z' \rangle = \langle z, x^* \rangle\) for all \(x^* \in \mathbf{E}^*\) whose class mod. \(^t u(\mathbf{F})\) is equal to \(z'\) ;
\item
  \(\langle t, t' \rangle = \langle s, t' \rangle\) for all \(s \in \mathbf{F}\) whose class mod. \(u(\mathbf{E})\) is equal to \(t\) .
\end{enumerate}

In particular we derive from these results:

PROPOSITION 10. Let E, F be two vector spaces over the same field K and \(u: \mathbf{E} \to \mathbf{F}\) a linear mapping.

\begin{enumerate}
\def\labelenumi{(\roman{enumi})}
\item
  For \(u\) to be injective (resp. surjective), it is necessary and sufficient that \(^{t}u\) be surjective (resp. injective).
\item
  \(\operatorname{rg}(u) \leqslant \operatorname{rg}(^t u)\) and \(\operatorname{rg}(u) = \operatorname{rg}(^t u)\) if \(\operatorname{rg}(u)\) is finite.
\end{enumerate}

The second assertion follows from the above and Theorem 4.

THEOREM 7. Let E be a vector space over a field K, F a subspace of E and F' the orthogonal of F in E*.

\begin{enumerate}
\def\labelenumi{(\roman{enumi})}
\item
  \(\dim \mathbf{F}' \geqslant \operatorname{codim}_{\mathbb{E}} \mathbf{F}\) ; for \(\dim \mathbf{F}'\) to be finite, it is necessary and sufficient that \(\operatorname{codim}_{\mathbb{E}} \mathbf{F}\) be finite, and then \(\dim \mathbf{F}' = \operatorname{codim}_{\mathbb{E}} \mathbf{F}\) .
\item
  The orthogonal of \(\mathbf{F}'\) in \(\mathbf{E}\) is equal to \(\mathbf{F}\) .
\item
  Every finite-dimensional subspace \(\mathbf{G}'\) of \(\mathbf{E}^*\) is the orthogonal of some subspace of \(\mathbf{E}\) , necessarily equal to the orthogonal of \(\mathbf{G}'\) in \(\mathbf{E}\) and of finite codimension.
\item
  We know that \(\mathbf{F}'\) is isomorphic to the dual \((\mathbf{E} / \mathbf{F})^*\) (§ 2, no. 6, Corollary to Proposition 9) and hence the assertion follows from Theorem 4, since \(\dim (\mathbf{E} / \mathbf{F}) = \operatorname{codim}_{\mathbf{E}}\mathbf{F}\) by definition.
\item
  Let \(F_{1}\) be the orthogonal of \(F'\) in E; clearly \(F \subset F_{1}\) and the orthogonal \(F_{1}'\) of \(F_{1}\) is equal to \(F'\) (§2, no. 4); the canonical linear mapping \((\mathrm{E}/\mathrm{F}_{1})^{*} \to (\mathrm{E}/\mathrm{F})^{*}\) , the transpose of \(E/F \to E/F_{1}\) is therefore bijective (§2, no. 6, Corollary to Proposition 10); it then follows from Proposition 10 that the canonical mapping \(E/F \to E/F_{1}\) is bijective, which implies \(F_{1} = F\) .
\item
  Let \(\mathbf{G}'\) be a subspace of \(\mathbf{E}^*\) of finite dimension \(p\) and let \(\mathbf{F}\) be its orthogonal in \(\mathbf{E}\) ; then \(\operatorname{codim}_{\mathbf{E}} \mathbf{F} \leqslant \dim \mathbf{G}'\) . For, if \((a_{i}^{*})_{1 \leqslant i \leqslant p}\) is a basis of \(\mathbf{G}'\) , \(\mathbf{F}\) is the kernel of the linear mapping \(x \mapsto (\langle x, a_{i}^{*} \rangle)\) from \(\mathbf{E}\) to \(\mathbf{K}_{s}^{p}\) whose rank is at most \(p\) (no. 4, Proposition 9), whence the conclusion (no. 4). Then let \(\mathbf{F}'\) be the orthogonal of \(\mathbf{F}\) in \(\mathbf{E}^*\) ; it follows from (i) that \(\dim \mathbf{F}' \leqslant \dim \mathbf{G}'\) ; but on the other hand obviously \(\mathbf{G}' \subset \mathbf{F}'\) , whence \(\mathbf{F}' = \mathbf{G}'\) ( \(\S 2\) , no. 3, Corollary 4 to Proposition 4).
\end{enumerate}

Remark. An infinite-dimensional subspace \(G'\) of \(E^{*}\) is not necessarily the orthogonal of a subspace of E, in other words, if F is the orthogonal of \(G'\) in E, the orthogonal \(F'\) of F in \(E^{*}\) may be distinct from \(G'\) (Exercise 20(b)).†

COROLLARY 1. Let \((x_{i}^{*})_{1\leqslant i\leqslant p}\) be a finite sequence of linear forms on E and let F be the subspace of E consisting of the \(x\) such that

\[
\langle x, x _ {i} ^ {*} \rangle = 0 \quad for 1 \leqslant i \leqslant p.
\]

Then \(\operatorname{codim}_{E} F\) is equal to the rank of the set of the \(x_{i}^{*}\) and every linear form on E which is zero on F is a linear combination of the \(x_{i}^{*}\) . Then \(\operatorname{codim}_{E} F \leqslant p\) and in order that \(\operatorname{codim}_{E} F = p\) , it is necessary and sufficient that the \(x_{i}^{*}\) be linearly independent.

The set \(G'\) of linear combinations of the \(x_{i}^{*}\) is a subspace of \(E^{*}\) and F is the orthogonal of \(G'\) in E, hence \(\operatorname{codim}_{E} F = \dim G'\) by Theorem 7; further \(\dim G' \leqslant p\) and the relation \(\dim G' = p\) means that \((x_{i}^{*})\) is a free system (no. 2, Proposition 1); hence the corollary.

COROLLARY 2. (i) Let \((x_{i}^{*})_{1\leqslant i\leqslant p}\) be a finite sequence of linear forms on E. For \((x_{i}^{*})\) to be a free system, it is necessary and sufficient that there exist a sequence \((x_{i})_{1\leqslant i\leqslant p}\) of elements of E such that \(\langle x_i,x_j^*\rangle = \delta_{ij}\) (Kronecker index).

\begin{enumerate}
\def\labelenumi{(\roman{enumi})}
\setcounter{enumi}{1}
\tightlist
\item
  Let \((x_{i})_{1\leqslant i\leqslant p}\) be a finite sequence of elements of E. For \((x_{i})\) to be a free system, it is necessary and sufficient that there exist a sequence \((x_{i}^{*})_{1\leqslant i\leqslant p}\) of linear forms on E such that \(\langle x_i,x_j^*\rangle = \delta_{ij}\) .
\end{enumerate}

Clearly (ii) follows from (i) by considering E as identified with a subspace of \(\mathbf{E}^{**}\) by means of \(c_{\mathbf{E}}\) (Theorem 6). Let \(\mathbf{G}'\) be the subspace of \(\mathbf{E}^*\) generated by the \(x_{i}^{*}\) and \(\mathbf{F}\) its orthogonal in \(\mathbf{E}\) ; \(\mathbf{E} / \mathbf{F}\) and \(\mathbf{G}'\) can each be canonically identified with the dual of the other; if the family \((x_{i}^{*})\) is free, there is in \(\mathbf{E} / \mathbf{F}\) a basis \((\dot{x}_i)\) the dual of \((x_{i}^{*})\) and every representative system \((x_{i})\) of the classes \(\dot{x}_{i}\) has the required properties. Conversely, the existence of the system \((x_{i})\) such that \(\langle x_{i}, x_{j}^{*} \rangle = \delta_{ij}\) implies that for all \(i\) the subspace of \(\mathbf{E}^*\) orthogonal to \(\mathbf{K}\) . \(x_{i}\) contains the \(x_{j}^{*}\) of index \(j \neq i\) but does not contain \(x_{i}^{*}\) , hence the system \((x_{i}^{*})_{1 \leqslant i \leqslant p}\) is free.

COROLLARY 3. Let S be a set and V a vector subspace of the right vector K-space \(K_{d}^{S}\) of mappings of S into K. In order that \(\dim V \geqslant p\) (where p is an integer), it is necessary and sufficient that there exist p elements \(s_{i}\) of S and p elements \(f_{i}\) of V ( \(1 \leqslant i \leqslant p\) ) such that \(f_{i}(s_{j}) = \delta_{ij}\) .

The space \(K_{d}^{S}\) is canonically identified with the dual of \(E = K_{d}^{(S)}\) and \(f(s) = \langle e_{s}, f \rangle\) for \(s \in S\) and \(f \in K_{s}^{S}\) , \((e_{s})_{s \in S}\) being the canonical basis of E. Corollary 2 then shows that the condition is sufficient. Conversely, suppose that \(\dim V \geqslant p\) , so that there exists a subspace \(G'\) of V of dimension p; let F be the orthogonal of \(G'\) in E, so that \(\dim(E/F) = p\) . It follows from no. 1, Theorem 2 that there exist p elements \(s_{i} \in S\) such that the \(e_{s_{i}} (1 \leqslant i \leqslant p)\) form a basis of a supplement F of E (applying Theorem 2 (no. 1) to a free subset consisting of a basis of E and the generating system the union of this free subset and the canonical basis of E); we then take the \(f_{i}\) to be the elements of a basis of G' dual to the basis of E/F consisting of the classes of the \(e_{s_{i}} \mod F\) .

COROLLARY 4. Let E be a vector space and M, N two subspaces of E of finite co-dimension; if M', N' are the orthogonals of M and N in E*, the orthogonal of M ∩ N in E* is M' + N'.

As M (resp. N) is the orthogonal of \(M'\) (resp. \(N'\) ) in E (Theorem 7), \(M \cap N\) is the orthogonal of \(M' + N'\) in E and hence \(M' + N'\) is the orthogonal of \(M \cap N\) in \(E^*\) (Theorem 7 (iii)).

COROLLARY 5. Let E be a vector space of finite dimension n.~For every subspace F of E of dimension p, the orthogonal F' of F in E* is of dimension n - p.~For every subspace G' of E* of dimension q, the orthogonal G of G' in E is of dimension n - q and G' is the orthogonal of G in E*.

Theorem 7 gives another characterization of hyperplanes in E:

PROPOSITION 11. For every hyperplane H in a vector space E, there exists a linear form \(x_0^*\) on E such that \(\mathrm{H} = x_0^{*-1}(0)\) . Given such a form \(x_0^*\) , for a linear form \(x^*\) on E to satisfy \(\mathrm{H} = x^{*-1}(0)\) , it is necessary and sufficient that \(x^* = x_0^*\alpha\) , where \(\alpha\) is a scalar \(\neq 0\) . Conversely, for every linear form \(x^* \neq 0\) on E, the subspace \(x^{*-1}(0)\) is a hyperplane of E.

This statement merely expresses Theorem 7 for subspaces of E of codimension 1 and subspaces of E* of dimension 1.

If H is a hyperplane and \(x_{0}^{*}\) a linear form such that \(\mathrm{H}=x^{*-1}(0)\) , the relation

\[
\langle x, x _ {0} ^ {*} \rangle = 0
\]

which characterizes the elements \(x \in \mathbf{H}\) , is called an equation of \(\mathbf{H}\) .

More generally, if \((x_{\iota}^{*})\) is a family of linear forms on E and F denotes the vector subspace the intersection of the hyperplanes \(x_{\iota}^{*-1}(0)\) , the relation ``for all \(\iota\) , \(\langle x, x_{\iota}^{*}\rangle = 0\) '' characterizes the elements of F; the relations

\[
\langle x, x _ {\iota} ^ {*} \rangle = 0 \quad \text { for   all } \iota
\]

form a system of equations of the subspace F. Theorem 7 (ii) expresses the fact that every vector subspace of E can be defined by a system of equations.

Theorem 7 (i) and (ii) proves moreover that a subspace F of finite codimension p can be defined by a system of p equations

\[
\langle x, x _ {i} ^ {*} \rangle = 0, \quad 1 \leqslant i \leqslant p,\tag{19}
\]

where the forms \(x_{i}^{*}\) are linearly independent. Conversely, Corollary 1 to Theorem 7 shows that a subspace \(F\) defined by a system of \(p\) equations (19) is of codimension \(\leqslant p\) and that it is of codimension \(p\) if and only if the \(x_{i}^{*}\) are linearly independent; it amounts to the same to say that F cannot be defined by a system consisting of at most p - 1 of the equations (19).

